% Preamble
\documentclass[11pt, aspectratio=169]{beamer}
	% Packages
	\usepackage[english]{babel}
	\usepackage{lipsum}
	\usepackage{mathptmx}
	\usepackage{xcolor}
	\usepackage{hyperref}
	\usepackage{url}
	\usepackage{multicol}
	% Themes
	\usefonttheme[onlymath]{serif}
	\usetheme[nofirafonts]{focus}
	\renewcommand{\thempfootnote}{\arabic{mpfootnote}}
	% Cover
	\title{Microeconometrics Using Stata}
	\subtitle{Stata Basics: Exercises}
	\author{
		Jhon R. Ordoñez
		\footnote{\label{1} National University of San Cristóbal de Huamanga}
		\footnote{\label{2} Faculty of Economic, Administrative and Accounting Sciences}
		\footnote{\label{3} Professional School of Economics}
	}
	\date{\today}
% Body
\begin{document}
	% Cover
	\begin{frame}
		\maketitle
	\end{frame}
	% Outline
	\begin{frame}[t]
		\frametitle{Outline}
		%\begin{multicols}{2}
			\tableofcontents
		%\end{multicols}
	\end{frame}
	% Section 1
	\section{Exercise 1}
		% Slide 1
		\begin{frame}
			\frametitle{Exercise 1}
				\begin{block}{Exercise 1}
					Find information on the estimation method \textcolor{blue}{\textbf{clogit}} using \textcolor{blue}{\textbf{help}} and
					\textcolor{blue}{\textbf{search}}. Comment on the relative usefulness of these search commands.
				\end{block}
		\end{frame}
	% Section 2
	\section{Exercise 2}
		% Slide 1
		\begin{frame}
			\frametitle{Exercise 2}
			\begin{block}{Exercise 2}
				Download the Stata example dataset \textcolor{blue}{\textbf{auto.dta}}. Obtain summary
				statistics for \textcolor{blue}{\textbf{mpg}} and \textcolor{blue}{\textbf{weight}} according 
				to whether the car type is foreign (use the \textcolor{blue}{\textbf{by foreign}}: prefix). Comment on any differences
				between foreign and domestic cars. Then, \textcolor{blue}{\textbf{regressmpg}} on \textcolor{blue}{\textbf{weight}} and
				\textcolor{blue}{\textbf{foreign}}. Comment on any difference for foreign cars.
			\end{block}
		\end{frame}
	% Section 3
	\section{Exercise 3}
		% Slide 1
		\begin{frame}
			\frametitle{Exercise 3}
			\begin{block}{Exercise 3}
				Write a do-file to repeat the previous question. This do-file should
				include a log file. Run the do-file, and then use a text editor to view the
				log file.
			\end{block}
		\end{frame}
	% Section 4
	\section{Exercise 4}
		% Slide 1
		\begin{frame}
			\frametitle{Exercise 4}
			\begin{block}{Exercise 4}
				Using \textcolor{blue}{\textbf{auto.dta}}, obtain summary statistics for the \textcolor{blue}{\textbf{price}} variable.
				Then, use the results stored in \textcolor{blue}{\textbf{r()}} to compute a scalar, \textcolor{blue}{\textbf{cv}}, equal to the
				coefficient of variation (the standard deviation divided by the mean) of \textcolor{blue}{\textbf{price}}.
			\end{block}
		\end{frame}
	% Section 5
	\section{Exercise 5}
		% Slide 1
		\begin{frame}
			\frametitle{Exercise 5}
			\begin{block}{Exercise 5}
				Using \textcolor{blue}{\textbf{auto.dta}}, \textcolor{blue}{\textbf{regress mpg}} on \textcolor{blue}{\textbf{price}} and \textcolor{blue}{\textbf{weight}}. Then, use the
				results stored in \textcolor{blue}{\textbf{e()}} to compute a scalar, \textcolor{blue}{\textbf{r2adj}}, equal to $\bar{R}^{2}$.
				The adjusted $R^{2}$ equals 
					\begin{equation*}
						\bar{R}^{2} = R^{2} - (1 - R^{2}) \left(\dfrac{K-1}{N - K}\right)
					\end{equation*}
				where $N$ is the number of observations and $K$ is the number of regressors including
				the intercept. Also, use the results stored in \textcolor{blue}{\textbf{e()}} to calculate a scalar,
				\textcolor{blue}{\textbf{tweight}}, equal to the $t$ statistic to test that the coefficient of \textcolor{blue}{\textbf{weight}} is zero.
			\end{block}
		\end{frame}
	% Section 6
	\section{Exercise 6}
		% Slide 1
		\begin{frame}
			\frametitle{Exercise 6}
			\begin{block}{Exercise 6}
				Using \textcolor{blue}{\textbf{auto.dta}}, define a global macro named \textcolor{blue}{\textbf{varlist}} for a variable
				list with \textcolor{blue}{\textbf{mpg}}, \textcolor{blue}{\textbf{price}}, and \textcolor{blue}{\textbf{weight}}, and then obtain summary statistics
				for \textcolor{blue}{\textbf{varlist}}. Repeat this exercise for a local macro named \textcolor{blue}{\textbf{varlist}}.
			\end{block}
		\end{frame}
	% Section 7
	\section{Exercise 7}
		% Slide 1
		\begin{frame}
			\frametitle{Exercise 7}
			\begin{block}{Exercise 7}
				Using \textcolor{blue}{\textbf{auto.dta}}, use a \textcolor{blue}{\textbf{foreach}} loop to create a variable, \textcolor{blue}{\textbf{total}}, equal
				to the sum of \textcolor{blue}{\textbf{headroom}} and \textcolor{blue}{\textbf{length}}. Confirm by using \textcolor{blue}{\textbf{summarize}} that
				\textcolor{blue}{\textbf{total}} has a mean equal to the sum of the means of \textcolor{blue}{\textbf{headroom}} and \textcolor{blue}{\textbf{length}}.
			\end{block}
		\end{frame}
	% Section 8
	\section{Exercise 8}
		% Slide 1
		\begin{frame}
			\frametitle{Exercise 8}
			\begin{block}{Exercise 8}
				Create a simulated dataset with $100$ observations on two random
				variables that are each drawn from the uniform distribution. Use a seed
				of $12345$. In theory, these random variables have a mean of $0.5$ and a
				variance of $1/12$. Does this appear to be the case here? 
			\end{block}
		\end{frame}
	% References
	\begin{frame}[t,allowframebreaks]
		\frametitle{References}
		\bibliography{biblio.bib}
		% Style
		\bibliographystyle{apalike}
		% Authors
		\nocite{cameron2022-I}
		\nocite{cameron2022-II}
	\end{frame}
	% Cover
	\begin{frame}
		\maketitle
	\end{frame}
\end{document}